\begin{conclusion}[临界点分裂域的带符号置换结构：\texorpdfstring{$C_{2}\wr S_{3}$}{C2 wr S3} 的显式呈现与判别式素数 \(37\)]\label{con:xi-terminal-zm-elliptic-lattes-critical-sextic-galois}
沿用结论 \ref{con:xi-terminal-zm-elliptic-lattes-doubling} 的记号，并记
\[
C(X):=2X^{6}-10X^{4}+10X^{3}-10X^{2}+2X+1\in\QQ[X],
\qquad
D(T):=4T^{3}-4T+1\in\QQ[T].
\]
设 \(\alpha_{1},\alpha_{2},\alpha_{3}\) 为 \(D(T)=0\) 的三根。
则在其分裂域中，临界六次多项式发生完全分解并具有闭式乘积结构
\[
\boxed{\;
C(X)=2\prod_{i=1}^{3}\bigl(X^{2}-2\alpha_{i}X-2\alpha_{i}^{2}+1\bigr)\; }.
\]
特别地，对每一 \(\alpha_i\)，临界点方程在 \(\Phi(X)=\alpha_i\) 的纤维上等价于二次方程
\[
X^{2}-2\alpha_{i}X-2\alpha_{i}^{2}+1=0,
\]
从而对应的两临界点为
\[
X_{i,\pm}=\alpha_{i}\pm\sqrt{\,3\alpha_{i}^{2}-1\,}.
\]
并且其范数满足严格恒等式
\[
\boxed{\;
\operatorname{N}_{\QQ(\alpha_{i})/\QQ}\!\bigl(3\alpha_{i}^{2}-1\bigr)=-\frac{37}{16}\; }.
\]
\par
此外，\(C(X)\) 在 \(\QQ\) 上不可约，其判别式为
\[
\boxed{\ \mathrm{Disc}(C)=-2^{8}\cdot 37^{5}=-2^{8}\cdot \Delta(E)^{5}\ }.
\]
其中 \(\Delta(E)=37\) 为椭圆曲线 \(E:\ Y^2=X^3-X+\tfrac14\) 的判别式。
因而临界点分裂域必含二次子域 \(\QQ(\sqrt{-37})\)，且除 \(2\) 与 \(37\) 外无其它素分歧。
更精确地，设 \(L=\mathrm{Spl}(C)\) 为 \(C\) 的分裂域，则其伽罗瓦群同构于三维带符号置换群（Coxeter 群 \(B_{3}\)）
\[
\boxed{\;
\mathrm{Gal}(L/\QQ)\ \cong\ C_{2}^{3}\rtimes S_{3}\ \cong\ C_{2}\wr S_{3}\ \cong\ S_{4}\times C_{2}\; } ,
\]
其阶为 \(48\)。
该结构可解释为：三次多项式 \(D(T)\) 的分裂域提供 \(S_{3}\)-对称的“临界值三重体”，而三组二次分支 \(\sqrt{3\alpha_{i}^{2}-1}\) 叠加给出 \(C_{2}^{3}\) 的符号翻转层；两者半直积给出全域对称性。
\end{conclusion}
\begin{proof}[证明要点]
若 \(\alpha\) 满足 \(D(\alpha)=0\)，则在 \(\QQ(\alpha)[X]\) 中存在严格恒等式
\[
\bigl(X^{4}+2X^{2}-2X+1\bigr)-\alpha\bigl(4X^{3}-4X+1\bigr)=\bigl(X^{2}-2\alpha X-2\alpha^{2}+1\bigr)^{2},
\]
由此得到所示二次因子与 \(C(X)\) 的三组分裂；并由二次判别式立即得到临界点闭式
\(
X_{i,\pm}=\alpha_i\pm\sqrt{3\alpha_i^{2}-1}
\)。
范数恒等式可由结果式计算
\(
\operatorname{N}_{\QQ(\alpha)/\QQ}(3\alpha^{2}-1)=\mathrm{Res}_{\alpha}\bigl(\alpha^{3}-\alpha+\tfrac14,\ 3\alpha^{2}-1\bigr)
\)
直接得到。
\par
不可约性、判别式分解与伽罗瓦群阶为直接符号计算结论；
其中 \(\mathrm{Gal}(L/\QQ)\cong S_{4}\times C_{2}\) 与 \(C_{2}^{3}\rtimes S_{3}\) 的同构为全八面体群的标准等价表述。
可复现核验脚本：\texttt{scripts/exp\_fold\_zm\_elliptic\_lattes\_rational\_points\_audit.py}。证毕。
\end{proof}

\begin{conclusion}[八面体 \texorpdfstring{$4$}{4}-挠射影场的 \texorpdfstring{$S_3$}{S3} 商：临界六次分裂域强制包含 \texorpdfstring{$E[2]$}{E[2]} 的伽罗瓦闭包]\label{con:xi-terminal-zm-elliptic-lattes-critical-sextic-s3-quotient}
沿用结论 \ref{con:xi-terminal-zm-elliptic-lattes-critical-sextic-galois} 的记号，设
\[
L:=\mathrm{Spl}\bigl(C(X)\bigr),\qquad L_{2}:=\mathrm{Spl}\bigl(D(T)\bigr).
\]
则有包含关系 \(L_2\subset L\)，并且由 \(C_2^3\rtimes S_3\) 的半直积结构给出规范满射
\[
\boxed{\;
\mathrm{Gal}(L/\QQ)\twoheadrightarrow \mathrm{Gal}(L_2/\QQ)\cong S_3\;}
\]
其核同构于 \(C_2^3\)（对应三组平方根 \(\sqrt{3\alpha_i^2-1}\) 的独立符号翻转）。
因此 \(L_2=\mathrm{Spl}\bigl(D_2(X)\bigr)\) 可识别为 \(\QQ(E[2])\) 的 \(S_3\)-伽罗瓦闭包（结论 \ref{con:xi-terminal-zm-collision-s2-e2-field-isomorphism}），从而 \(L\) 在群论意义上不可避免地支配 \(E[2]\) 的 \(S_3\)-闭包。
特别地，由结论 \ref{con:xi-terminal-zm-collision-s2-e2-field-isomorphism}，
\(
\QQ(\Lambda)\cong \QQ(E[2])
\)
为该 \(S_3\)-闭包的三次子域，因而亦嵌入 \(L\)。
\end{conclusion}
\begin{proof}[证明要点]
由结论 \ref{con:xi-terminal-zm-elliptic-lattes-critical-sextic-galois} 的分解式，
\(
C(X)=2\prod_{i=1}^{3}(X^{2}-2\alpha_{i}X-2\alpha_{i}^{2}+1)
\)，
可知 \(L\) 由 \(D(T)\) 的分裂域与三组二次分支生成：
\[
L=L_2\bigl(\sqrt{3\alpha_1^2-1},\sqrt{3\alpha_2^2-1},\sqrt{3\alpha_3^2-1}\bigr).
\]
对每个 \(i\)，平方根的符号翻转给出一个阶 \(2\) 自同构；三者生成的正规子群同构于 \(C_2^3\)，并在商上诱导对三元组 \((\alpha_1,\alpha_2,\alpha_3)\) 的置换作用。
因此存在短正合列
\(
1\to C_2^3\to \mathrm{Gal}(L/\QQ)\to \mathrm{Gal}(L_2/\QQ)\to 1
\)，
得到所示满射。
最后，\(D(T)\) 与结论 \ref{con:xi-terminal-zm-collision-s2-e2-field-isomorphism} 的 \(D_2(X)\) 同型（仅变量记号不同），其分裂域即为 \(E[2]\) 的 \(S_3\)-伽罗瓦闭包。
由该结论的三次域同构 \(\QQ(\Lambda)\cong\QQ(E[2])\)，立得嵌入断言。证毕。
\end{proof}

